\documentclass[11pt]{article}
\usepackage{fontspec}
\usepackage[french,english]{babel}
\usepackage[leqno]{amsmath}
\usepackage{amssymb}
\usepackage{amsthm}
\usepackage{mathrsfs}
\usepackage{tikz-cd}
\usepackage{enumitem}

\newcommand{\printedpage}[1]{\marginpar{\footnotesize p.~#1}}
\newcommand{\msp}[1]{\marginpar{\footnotesize [#1]}}

% Printed arrow notations (glyphs approximated):
% Q-orthogonal couple: harpoon arrow with hooked/barbed ends
\newcommand{\Qorth}{\mathrel{\leftharpoondown\mkern-6mu\relbar\mkern-6mu\rightharpoonup}}
% the same arrow crossed by a bar
\newcommand{\Qorthx}{\mathrel{\leftharpoondown\mkern-9mu{+}\mkern-9mu\rightharpoonup}}
% harpoon tail, cross, full arrow head (as printed in a) of the first page)
\newcommand{\Qorthxfull}{\mathrel{\leftharpoondown\mkern-9mu{+}\mkern-9mu\rightarrow}}
% couple de Quillen: double arrow crossed by a bar
\newcommand{\Qcouple}{\mathrel{\leftarrow\mkern-9mu{+}\mkern-9mu\rightarrow}}
% closed couple: long double arrow
\newcommand{\clos}{\longleftrightarrow}
% long arrow with small head
\newcommand{\fl}{\longrightarrow}

\begin{document}

\printedpage{4}
% head: A. Grothendieck : Dérivateurs, Ch. XIX, Édition des Universités de Montpellier II et Paris VII
% foot: Version 11 février 2019

\msp{page 6}
\textit{then}
\begin{enumerate}[label={}]
\item[a)] \textit{one has $\Phi_1 \Qorthxfull \Psi$ and} % reading resolved: DER_XIX_corrections.csv
\item[b)] $\Phi_1^{\sim} = \Phi^{\sim} = \Phi_1^{\natural}$.
\end{enumerate}
This corollary follows from the theorem by virtue of XV prop. 4.2, of which I recall here the essentials:

\textbf{Proposition 1.3.} \textit{Let $\Phi \Qorthx \Psi$. Then $\Phi^{\sim} = \Phi^{\natural} = \Psi^{*}$, and dually one has $\Psi^{-} = \Psi^{\natural} = \Phi_{*}$.}

Let us also note, in this line of thought:

\textbf{Corollary 1.4.} \textit{Let $\Phi \Qorth \Psi$, and let $W$ be an arbitrary subset of $\mathrm{Fl}(\mathcal{M})$. In order to have $\Phi \subset W^{\natural}$, it suffices that every arrow of $\mathcal{M}$ factor as $f = pi$, with $i \in W$, $p \in \Psi$. (And this condition is obviously also necessary if $\Phi \Qorthx \Psi$, and $W$ stable under direct factors.)}

One says that a subset $\Phi_0$ of $\Phi$ is Q-\textit{generating on the left} for $\Phi$ if it satisfies the condition
\[
\Phi \ \subset \Phi_0^{\sim}\ , \tag*{(1.9)}
\]
and one says that $\Phi$ is a \textit{special} subset if $\Phi$ admits a \textit{small} left Q-generating subset. If $\Phi_0 \subset \Phi$ is small, and if $\mathcal{M}$ satisfies the conditions of theorem 1.1, then $\Phi_0$ is a Q-generating subset (on the left) if and only if its total Q-saturation (on the left) $\Phi_0'''$ contains $\Phi$ (a condition apparently stronger than $\Phi \subset \Phi_0^{\sim}$, since in general one has $\Phi_0''' \subset \Phi_0^{\sim}$). This therefore also means, via 1.2, that every \msp{page 7} element of $\Phi$ is a direct factor of a transfinite composite of arrows obtained from
\[
\Phi_0 \cup (\text{Identities of } \mathcal{M})
\]
by cobase change.

I recall also the $\pm$ tautological sorites of adjoint functors:

\textbf{Proposition 1.5.} \textit{Let}
\[
\begin{tikzcd}
\mathcal{M}' \arrow[r, shift right=1, "v^{*}"'] & \mathcal{M} \arrow[l, shift right=1, "v_{!}"']
\end{tikzcd}
\tag*{(1.5.1)}
\]
\textit{be a pair of adjoint functors, let}
\[
\Phi \ \subset\ \mathrm{Fl}(\mathcal{M})\ , \qquad \Psi' \ \subset\ \mathrm{Fl}(\mathcal{M}')\ ,
\]
\textit{then one has}
\[
(v_{!}(\Phi) \Qorth \Psi') \qquad \Longleftrightarrow \qquad (\Phi \Qorth v^{*}(\Psi'))\ . \tag*{(1.5.2)}
\]

\textbf{Corollary 1.6.} \textit{For every subset $\Phi$ of} $\mathrm{Fl}(\mathcal{M})$,
\[
(v_{!}(\Phi))_{*} \ =\ (v^{*})^{-1}(\Phi_{*})\ . \tag*{(1.6.1)}
\]

\printedpage{5}
% head: A. Grothendieck : Dérivateurs, Ch. XIX, Édition des Universités de Montpellier II et Paris VII
% foot: Version 11 février 2019

Indeed, applying 1.5 to $\Phi$ and to $\{f'\}$, where $f' \in \mathrm{Fl}(\mathcal{M}')$, one finds
\[
f' \in (v_{!}(\Phi))_{*} \qquad \Longleftrightarrow \qquad v^{*}(f') \in \Phi_{*}\ ,
\]
which is equivalent to (1.6.1). Dually:

\textbf{Corollary 1.7.} \textit{For every subset $\Psi'$ of $\mathcal{M}'$, one has}
\[
(v^{*}(\Psi'))^{*} \ =\ (v_{!})^{-1}(\Psi'^{*}). \tag*{(1.7.1)}
\]

Indeed, if $f \in \mathrm{Fl}(\mathcal{M})$, one has by (1.5.2)
\[
f \in (v^{*}(\Psi'))^{*} \qquad \Longleftrightarrow \qquad v_{!}(f) \in \Psi^{*}\ . % [sic] \Psi^{*} -> \Psi'^{*}, cf. (1.7.1)
\]
\textsuperscript{(}\footnote[3]{give here 1.9 page 9.}\textsuperscript{)}.

\msp{page 8}
A Q-orthogonal couple $\Phi \Qorth \Psi$ is said to be \textit{special} if it satisfies the following conditions.
\[
\left\{
\begin{array}{ll}
\text{a)} & \parbox[t]{0.7\textwidth}{$\mathcal{M}$ is stable under small $\varinjlim$ , the objects of $\mathcal{M}$ are accessible. (One then says that $\mathcal{M}$ is ``special''.)} \\[3ex]
\text{b)} & \text{The couple is closed: } \Phi \clos \Psi. \\[1ex]
\text{c)} & \Phi \text{ is \textit{special.}}
\end{array}
\right.
\tag*{(1.10)}
\]
This therefore implies, by 1.1, that $\Phi \Qcouple \Psi$ is a Quillen couple. If $\mathcal{M}$ satisfies a), the special Q-orthogonal couples in $\mathcal{M}$, or again the special Quillen couples, thus correspond one-to-one to the subsets $\Phi$ of $\mathrm{Fl}(\mathcal{M})$ which are \textit{special} and left closed (or what amounts to the same thing, left Q-saturated), by associating to such a $\Phi$ the couple $\Phi \Qcouple \Phi_{*}$.

This being so:

\marginpar{$\Vert$}% double vertical bar printed in the left margin beside Corollaire 1.8
\textbf{Corollary 1.8.} \textit{Suppose that $\mathcal{M}$, $\mathcal{M}'$ satisfy condition} a) \textit{of} (1.10), \textit{and let $(\Phi, \Psi)$ be a special Quillen couple in $\mathcal{M}$. Then $\Bigl(\Phi' = v_{!}(\Phi)^{\sim}, \Psi' = (v^{*})^{-1}(\Psi)\Bigr)$ is a special Quillen couple in $\mathcal{M}'$.}

Indeed, by corollary 1.6 one has
\[
\Phi'_{*} \ =\ (v_{!}(\Phi))_{*} \ =\ (v^{*})^{-1}(\Psi) \ =\ \Psi'\ ,
\]
hence $\Phi' \clos \Psi'$. It remains to prove that $\Phi'$ is special. Let $\Phi_0$ be a small subset of $\Phi$ such that $\Phi \subset \Phi_0^{\sim}$, put $\Phi'_0 \ =\ v_{!}(\Phi_0) \subset \Phi'$, \msp{page 9} it suffices to see that one has
\[
v_{!}(\Phi) \ \subset\ v_{!}(\Phi_0)^{\sim}\ , \tag*{($*$)}
\]
from which it will follow that one has $\Phi' \overset{\text{def}}{=} (v_{!}(\Phi))^{\sim} \subset v_{!}(\Phi_0)^{\sim}$. But one has
\[
v_{!}(\Phi_0)^{\sim} \ =\ (v_{!}(\Phi_0)_{*})^{*}\ ,
\]
\printedpage{6}
% head: A. Grothendieck : Dérivateurs, Ch. XIX, Édition des Universités de Montpellier II et Paris VII
% foot: Version 11 février 2019
and
\[
v_{!}(\Phi_0)_{*} \ =\ (v^{*})^{-1}(\underbrace{\Phi_{0\,*}}_{\Psi}) \ =\ (v^{*})^{-1}(\Psi) \ \overset{\text{def}}{=}\ \Psi'\ ,
\]
hence
\[
v_{!}(\Phi_0)^{\sim} \ =\ \Psi'^{*} \ =\ \Phi' \ =\ v_{!}(\Phi)^{\sim}\ ,
\]
whence ($*$). This argument does not appeal to the ``special'' hypothesis on $\mathcal{M}$, and can be made explicit in the

\textbf{Corollary 1.9.} \textsuperscript{(}\footnote[4]{should come before 1.8.}\textsuperscript{)} \textit{Under the general conditions of} 1.5 \textit{($v_{!} : \mathcal{M} \fl \mathcal{M}'$ being a functor admitting a right adjoint), let $\Phi_0$, $\Phi$ be two subsets of} $\mathrm{Fl}(\mathcal{M})$ \textit{such that $\Phi \subset \Phi_0^{\sim}$ (or what amounts to the same thing, $\Phi^{\sim} \subset \Phi_0^{\sim}$). Then $v_{!}(\Phi) \subset v_{!}(\Phi_0)^{\sim}$, i.e. $v_{!}(\Phi)^{\sim} \subset v_{!}(\Phi_0)^{\sim}$.}

Indeed, this relation is also equivalent to $v_{!}(\Phi_0)_{*} \ \subset\ v_{!}(\Phi)_{*}$, that is (by 1.6) $(v^{*})^{-1}(\Phi_{0\,*}) \ \subset\ (v^{*})^{-1}(\Phi_{*})$, and is implied by $\Phi_{0\,*} \ \subset\ \Phi_{*}$. The latter on the other hand is equivalent to the hypothesis $\Phi^{\sim} \subset \Phi_0^{\sim}$, q.e.d.

\msp{page 10}
If $\Psi$ is a subset of $\mathrm{Fl}(\mathcal{M})$, one obtains from it for every $I$ in Cat a subset
\[
\left\{
\begin{array}{l}
\Psi(I) \subset \mathrm{Fl}(\mathcal{M}(I)), \quad \text{with}\ \ \mathcal{M}(I) = \underline{\mathrm{Hom}}(I^{\circ}, \mathcal{M})\ , \\[1ex]
f \in \Psi(I) \iff f(i) \in \Psi \ \ \forall\, i \in \mathrm{Ob}\, I\ .
\end{array}
\right.
\tag*{(1.11)}
\]
If one introduces the discrete category $I_0$ underlying $\mathrm{Ob}\, I$, and the canonical functor
\[
v_I = v \ :\ I_0 \fl I\ , \tag*{(1.12)}
\]
one finds a pair of adjoint functors
\[
\begin{tikzcd}[row sep=small]
\mathcal{M}(I) \arrow[r, shift right=1, "v^{*}"'] & \mathcal{M}(I_0) \arrow[l, shift right=1, "v_{!}"'] \\
( & ) \arrow[l, dotted, "v_{*}"]
\end{tikzcd}
\ \simeq\ \mathcal{M}^{I_0}\ ,
\tag*{(1.13)}
\]
and one has
\[
\Psi(I) \ =\ (v^{*})^{-1}(\Psi(I_0))\ .
\]
Suppose now that one has
\[
\left\{
\begin{array}{rcl}
\Psi & = & \Phi_{*} \\[1ex]
\text{whence}\ \ \Psi(I_0) & = & \Phi(I_0)_{*}
\end{array}
\right.
\tag*{(1.14)}
\]
(immediate verification), one then concludes by 1.6
\[
\Psi(I) \ =\ (v_{!}(\Phi(I_0)))_{*} \ =\ (\Phi\langle I\rangle)_{*}\ , \tag*{(1.15)}
\]
with
\[
\Phi\langle I\rangle \ \overset{\text{def}}{=}\ v_{!}(\Phi(I_0))^{\sim} \ =\ \Psi(I)^{*}\ . \tag*{(1.16)}
\]
If one has a functor
\[
u \ :\ I \fl I'\ , \tag*{(1.17)}
\]
whence a triple of adjoint functors
\printedpage{7}
% head: A. Grothendieck : Dérivateurs, Ch. XIX, Édition des Universités de Montpellier II et Paris VII
% foot: Version 11 février 2019
\msp{page 11}
\[
\left\{
\begin{tikzcd}[column sep=large]
\mathcal{M}(I) \arrow[r, shift left=3, "u_{!}"] \arrow[r, shift right=3, "u_{*}"] & \mathcal{M}(I') \arrow[l, "u^{*}"']
\end{tikzcd}
\right.
,
\tag*{(1.18)}
\]
one will have in a $\pm$ tautological way
\[
\left\{
\begin{array}{rcl}
u_{!}(\Phi\langle I\rangle) & \subset & \Phi\langle I'\rangle \\
u^{*}(\Psi(I')) & \subset & \Psi(I)\ .
\end{array}
\right.
\tag*{(1.19)}
\]
Corollary 1.8 gives:

\textbf{Proposition 1.10.} \textit{Suppose that $\Phi \Qcouple \Psi$ is a} special \textit{Quillen couple in $\mathcal{M}$, i.e. of the form $(\Phi, \Phi_{*})$, with $\Phi$ a special left Q-closed subset in $\mathcal{M}$ ($\mathcal{M}$ satisfying condition} (1.10. a))\textit{. Then for every $I$ in} Cat, \textit{$(\Phi\langle I\rangle, \Psi(I))$ is a special couple in $\mathcal{M}(I)$, in particular it satisfies the factorization theorem.}

Let us consider the dual construction, which leads us to introduce, starting from a couple $\Phi \clos \Psi$ in $\mathcal{M}$, the couple
\[
\Phi(I) \ \clos\ \Psi[I]\ , \tag*{(1.20)}
\]
where $\Phi(I)$ is defined by the general formula (1.11) (where $\Psi$ is replaced by $\Phi$), and where $\Psi[I]$ is defined either as $\Phi(I)_{*}$, or as
\[
\Psi[I] \ =\ v_{*}(\Psi(I_0))^{-} \ =\ \Phi(I)_{*}\ . \tag*{(1.21)}
\]

\msp{page 12}
One then has, for $u : I \fl I'$
\[
\left\{
\begin{array}{rcl}
u^{*}(\Phi(I')) & \subset & \Phi(I) \\
u_{*}(\Psi[I]) & \subset & \Psi[I']\ .
\end{array}
\right.
\tag*{(1.22)}
\]
There still arises the question whether the couple $(\Phi(I), \Psi[I])$ satisfies the factorization theorem, under suitable hypotheses on $(\Phi, \Psi)$. I can only venture the

\textbf{Question-conjecture 1.11.} \textit{Suppose that $(\Phi, \Psi)$ is a} special \textit{Quillen couple, is the same true of $(\Phi(I), \Psi[I])$ for every $I$ in} (Cat)\textit{? In other words, if $\mathcal{M}$ is special, condition} (1.10. a))\textit{, and if $\Phi$ is a} special \textit{left Q-closed subset of $\mathcal{M}$, i.e. admitting a small Q-generating subset, is it true that for every $I$ in} Cat, \textit{$\Phi(I)$ is also special? I conjecture that yes, making if need be additional hypotheses on $\mathcal{M}$ (existence of a cardinal filtration?). E.g. if $\Phi_0$ is Q-generating in $\Phi$, and small, then is $\Phi_0(I)$ Q-generating in $\Phi(I)$?}

\textbf{Proposition 1.12} \textsuperscript{(}\footnote[5]{cf. XVII 6.4 (p. 117), but where the notations for $\Phi\langle I\rangle$, $\Psi(I)$ are reversed.}\textsuperscript{)}. \textit{Let $\Phi \clos \Psi$ in $\mathcal{M}$, let $u : I \fl I'$ be an arrow in} Cat, \textit{whence}
\[
\left\{
\begin{tikzcd}[column sep=large]
\mathcal{M}(I) \arrow[r, shift left=3, "u_{!}"] \arrow[r, shift right=3, "u_{*}"] & \mathcal{M}(I') \arrow[l, "u^{*}"']
\end{tikzcd}
\right.
.
\tag*{(1.12.1)}
\]
\printedpage{8}
% head: A. Grothendieck : Dérivateurs, Ch. XIX, Édition des Universités de Montpellier II et Paris VII
% foot: Version 11 février 2019
\textit{One then has} (1.19), \textit{i.e.}
\[
\left\{
\begin{array}{rcl}
u_{!}(\Phi\langle I\rangle) & \subset & \Phi\langle I'\rangle \\
u^{*}(\Psi(I')) & \subset & \Psi(I)\ ,
\end{array}
\right.
\tag*{(1.13.1)}
\]
\textit{and moreover the equivalence}
\[
u^{*}(\Phi(I')) \subset \Phi\langle I'\rangle \qquad \Longleftrightarrow \qquad u_{*}(\Psi(I)) \subset \Psi(I')\ . \tag*{(1.13.2)}
\]
\textit{These equivalent relations are satisfied in each of the three cases}~\textsuperscript{(}\footnote[6]{\textbf{N.B.} One has $1^{\circ} \Longrightarrow 3^{\circ}$, $2^{\circ} \Longrightarrow 3^{\circ}$.}\textsuperscript{)} \textit{:}
\begin{enumerate}[label={}]
\item[1°)] \textit{$u$ is coétale.}
\item[2°)] \textit{$I$ is discrete.}
\item[3°)] $\forall\, i' \in \mathrm{Ob}\, I'$\textit{, every connected component of $I/i'$ admits a final object.}
\end{enumerate}
\msp{page 13}

\section*{2\quad Reminders on closed Quillen triples and special Quillen triples}

Let $(\mathcal{M}, W)$ be a \textit{category of models,} i.e. $\mathcal{M}$ equipped with a localizer $W \subset \mathrm{Fl}(\mathcal{M})$, i.e. a subset satisfying
\[
\parbox{0.75\textwidth}{$W$ contains the isomorphisms, and if $f$, $g$ are composable arrows, if two among $f$, $g$, $gf$ are in $W$, the third also \textit{(weak saturation).}}
\tag*{(2.1)}
\]
I recall that the datum of a \textit{closed Quillen triple} $(W, C, F)$\textit{,} or of a closed Quillen structure, on $\mathcal{M}$, having $W$ as set of quasi-isomorphisms, amounts to the datum of two Quillen couples
\[
C \Qcouple TF\ , \qquad TC \ \Qcouple\ F\ , \tag*{(2.2)}
\]
satisfying
\[
TC \ \subset\ C \cap W\ , \qquad TF \ \subset\ F \cap W\ , \tag*{(2.3)}
\]
plus one of the three equivalent conditions:
\[
\left\{
\begin{array}{lrcll}
\text{a)} & TC & = & C \cap W & \qquad (\text{i.e. } C \cap W \subset TC) \\[1ex]
\text{b)} & TF & = & F \cap W & \qquad (\text{i.e. } F \cap W \subset TF) \\[1ex]
\text{c)} & W & = & TF \circ TC\ . &
\end{array}
\right.
\tag*{(2.4)}
\]
Let us also recall that this implies that $W$ \textit{is strongly saturated,} and a fortiori that $W = W^{\natural}$, i.e. $W$ stable under direct factors.

Let us also recall that by virtue of (1.4), one has:

\end{document}
